\documentclass[letterpaper,twocolumn,10pt]{article}
\usepackage{usenix-2020-09}

\begin{document}

\date{}
\title{\Large \bf Related Work}
\author{{\rm Jordan Guimfack Jeuna}\\Hochschule Bremerhaven}
\maketitle

\section{Related Work}
\label{sec:related}

Die Erkennung von Cyberangriffen in IoT-Netzwerken mittels Machine
Learning ist ein aktives Forschungsfeld, in dem realistische Benchmarks,
robuste Bewertungsmetriken und Erklärbarkeit zunehmend an Bedeutung
gewinnen. Die folgende Darstellung gliedert die relevante Literatur in
vier thematische Schwerpunkte.

\subsection{Datensatz und Benchmark}

Die empirische Grundlage dieser Arbeit bildet der CICIoT2023-Datensatz
von Neto et al.~\cite{neto2023ciciot2023}, der Netzwerkverkehr von 105
realen IoT-Geräten unter 33 Angriffsbedingungen in sieben Kategorien
erfasst: DDoS, DoS, Mirai, Reconnaissance, Spoofing, Brute-Force und
Web-based Attacks. Da alle Angriffe ausschließlich von IoT-Geräten gegen
IoT-Geräte ausgeführt wurden, bildet der Datensatz reale
Bedrohungsszenarien besonders realitätsnah ab, was ihn gegenüber
synthetischen Alternativen auszeichnet. Gleichzeitig weist er ein
ausgeprägtes Klassenungleichgewicht auf, das methodische Sorgfalt bei
der Modellbewertung erfordert.

\subsection{ML-basierte Angriffserkennung und Klassenungleichgewicht}

Surveys von Hamedani et al.~\cite{hamedani2025iotsurvey} sowie
Kikissagbe und Adda~\cite{kikissagbe2024review} belegen die breite
Nutzung ML-basierter IDS im IoT und heben baumbasierte
Ensemble-Methoden als konsistent leistungsstärkste Klasse hervor,
während Interpretierbarkeit und Ressourceneffizienz als weiterhin offene
Herausforderungen identifiziert werden.

Auf CICIoT2023 evaluieren Raturi et al.~\cite{raturi2026exploratory}
mehrere ML-Algorithmen und zeigen, dass Gradient Boosting in der
Mehrklassenklassifikation die höchste Balanced Accuracy erzielt. Als
konzeptioneller Ausgangspunkt dieser Arbeit benennen sie drei
methodische Lücken: fehlende Erklärbarkeit, fehlende
Fehlerfortpflanzungsanalyse und fehlende Latenzmessungen. Die
vorliegende Arbeit adressiert die erste dieser Lücken eigenständig
durch eine vollständige SHAP-Integration. Almahaqeri et
al.~\cite{almahaqeri2026gradient} evaluieren Gradient Boosting auf
demselben Datensatz mit stratifiziertem Undersampling als alternative
Imbalance-Strategie; Alharby~\cite{alharby2025ciciot} setzt zusätzlich
PCA als Dimensionsreduktionstechnik ein und liefert durch dokumentierte
Trainings- und Vorhersagezeiten realistische Vergleichswerte für die
eigene Ablation Study.

Hosseini et al.~\cite{hosseini2025imbalance} zeigen systematisch, dass
einfache Accuracy bei unausgewogenen Datensätzen zu verzerrten Urteilen
führt, da Modelle, die seltene Angriffsklassen ignorieren, rechnerisch
hohe Werte erzielen können. Raturi et al.~\cite{raturi2026exploratory}
belegen dies empirisch: ein Hidden Markov Model erreicht zwar
Recall~0{,}999 für DoS-Angriffe, lässt dabei jedoch 87\,\% aller
Nicht-DoS-Angriffe unerkannt. Sie schlagen daher Balanced Accuracy als
fairere Bewertungsgrundlage vor, einen Ansatz, den die vorliegende
Arbeit als primäre Metrik übernimmt.

\subsection{Erklärbarkeit und Forschungslücke}

Ein PRISMA-basierter Review von 129 Studien durch Ogunseyi et
al.~\cite{ogunseyi2026xai_review} belegt, dass SHAP und LIME in
98\,\% aller XAI-IDS-Arbeiten eingesetzt werden, ohne die
Erkennungsgenauigkeit zu beeinträchtigen, was die SHAP-Integration der
vorliegenden Arbeit wissenschaftlich begründet. Mohale und
Obagbuwa~\cite{mohale2025xai} belegen auf einem IDS-Datensatz, dass
SHAP konsistentere Erklärungen als LIME und ELI5 liefert und verborgene
Schwächen aufdeckt, die reine Accuracy-Metriken verbergen. Ihre Arbeit
dient als direkte methodische Vorlage für die eigene SHAP-Analyse.
Hermosilla et al.~\cite{hermosilla2025xai_forensic} quantifizieren die
Erklärungsqualität mit den Metriken Fidelität, Konsistenz und
Jaccard-Ähnlichkeit auf XGBoost- und TabNet-Modellen und zeigen, dass
SHAP LIME in Fidelität und Konsistenz übertrifft. Diese Kriterien
fundieren die semantische Validierung der SHAP-Ergebnisse in der
vorliegenden Arbeit.

Die Analyse der Literatur offenbart eine direkt belegbare
Forschungslücke: Keine der identifizierten Studien vereint auf
CICIoT2023 gleichzeitig Balanced Accuracy als primäre Metrik, eine
SHAP-Integration und eine Ablation Study zum Einfluss der
PCA-Reduktion. Raturi et al.~\cite{raturi2026exploratory} verwenden
Balanced Accuracy auf CICIoT2023, verzichten aber auf jede Form der
Erklärbarkeitsanalyse. Almahaqeri et al.~\cite{almahaqeri2026gradient}
und Alharby~\cite{alharby2025ciciot} integrieren weder Balanced
Accuracy als primäre Metrik noch eine SHAP-Komponente. Mohale und
Obagbuwa~\cite{mohale2025xai} sowie Hermosilla et
al.~\cite{hermosilla2025xai_forensic} liefern fundierte SHAP-Analysen,
arbeiten jedoch nicht auf CICIoT2023. Die vorliegende Arbeit schließt
diese Lücke, indem sie einen Gradient-Boosting-Klassifikator auf
CICIoT2023 mit Balanced Accuracy und einer vollständigen SHAP-Analyse
auf beiden Klassifikationsstufen verbindet.

\begin{thebibliography}{10}

\bibitem{neto2023ciciot2023}
E.~C.~P. Neto, S.~Dadkhah, R.~Ferreira, A.~Zohourian, R.~Lu, and
  A.~A. Ghorbani.
\newblock CICIoT2023: A Real-Time Dataset and Benchmark for Large-Scale
  Attacks in IoT Environment.
\newblock \emph{Sensors}, 23(13):5941, 2023. doi:10.3390/s23135941.
  \url{https://www.mdpi.com/1424-8220/23/13/5941}.
\newblock \textit{Wissenschaftliche Grundlage der Datenbeschreibung.
  Beantwortet die Frage: Wie wurde der CICIoT2023-Datensatz aufgebaut,
  welche IoT-Geräte und Angriffstypen sind enthalten, und warum ist er
  ein valider Benchmark?}

\bibitem{hamedani2025iotsurvey}
P.~Hamedani et~al.
\newblock Machine Learning-Based Security Solutions for IoT Networks:
  A Comprehensive Survey.
\newblock \emph{Sensors}, 25(11):3341, 2025. doi:10.3390/s25113341.
  \url{https://www.mdpi.com/1424-8220/25/11/3341}.
\newblock \textit{Survey über ML-basierte Sicherheitslösungen für IoT
  (2020--2024). Beantwortet die Frage: Warum ist ML-basierte
  Angriffserkennung für IoT ein relevantes und aktives Forschungsfeld,
  und welche methodischen Lücken bestehen noch? Wird im Theorieteil zur
  Motivation des Forschungsthemas verwendet.}

\bibitem{kikissagbe2024review}
B.~R. Kikissagbe and M.~Adda.
\newblock Machine Learning-Based Intrusion Detection Methods in IoT
  Systems: A Comprehensive Review.
\newblock \emph{Electronics}, 13(18):3601, 2024.
  doi:10.3390/electronics13183601.
  \url{https://www.mdpi.com/2079-9292/13/18/3601}.
\newblock \textit{Systematischer Review über supervised, unsupervised
  und hybride ML-Ansätze für IoT-IDS. Beantwortet die Frage: Welchen
  Platz nehmen baumbasierte Ensemble-Methoden im aktuellen
  Forschungsstand zu ML-basierten IDS für IoT ein, und wie wird die
  eigene Modellwahl theoretisch verankert?}

\bibitem{raturi2026exploratory}
A.~Raturi, R.~Aryan, and P.~Prashant.
\newblock An Exploratory Study on Application of Machine Learning in
  Detecting Cyberattacks.
\newblock \emph{Security and Privacy}, 9:e70220, 2026.
  doi:10.1002/spy2.70220.
  \url{https://onlinelibrary.wiley.com/doi/10.1002/spy2.70220}.
\newblock \textit{Ausgangspaper der Arbeit. Konzeptioneller Startpunkt,
  kein Reproduktionsziel. Benennt drei methodische Lücken: fehlende
  Erklärbarkeit, fehlende Fehlerfortpflanzungsanalyse, fehlende
  Latenzmessungen. Die eigene Arbeit adressiert Lücke 1 eigenständig
  durch SHAP-Integration.}

\bibitem{almahaqeri2026gradient}
S.~A. Almahaqeri, M.~H. Almourish, A.~A. Nasser, A.~A.~K. Elsayed,
  and A.~N. Alhejoj.
\newblock An Optimized Gradient Boosting Framework for IoT Intrusion
  Detection: A Comprehensive Evaluation on the CICIoT2023 Dataset.
\newblock \emph{Scientific Reports}, 2026. doi:10.1038/s41598-026-47399-5.
  \url{https://www.nature.com/articles/s41598-026-47399-5}.
\newblock \textit{Gleicher Datensatz (CICIoT2023), gleiche
  Modellkategorie (Gradient Boosting), andere Strategie für
  Klassenungleichgewicht (stratifiziertes Undersampling). Beantwortet
  die Frage: Wie unterscheidet sich die eigene hierarchische
  Klassifikation von alternativen Imbalance-Strategien auf demselben
  Datensatz?}

\bibitem{alharby2025ciciot}
M.~Alharby.
\newblock Evaluating Machine Learning Approaches for Multiple Attack
  Classification with Improved Computational Efficiency in IoT Networks.
\newblock \emph{Scientific Reports}, 15:39914, 2025.
  doi:10.1038/s41598-025-23711-7.
  \url{https://www.nature.com/articles/s41598-025-23711-7}.
\newblock \textit{Evaluiert Gradient Boosting und Decision Tree auf
  CICIoT2023 mit PCA als Dimensionsreduktion und misst Trainings- sowie
  Vorhersagezeiten. Beantwortet die Frage: Wie wirkt sich PCA auf die
  Leistung baumbasierter Modelle auf dem CICIoT2023-Datensatz aus, und
  was sind realistische Vergleichswerte für die eigene Ablation Study?}

\bibitem{hosseini2025imbalance}
S.~S. Hosseini et~al.
\newblock Addressing Class Imbalance in Intrusion Detection: A
  Comprehensive Evaluation of Machine Learning Approaches.
\newblock \emph{Electronics}, 14(1):69, 2025.
  doi:10.3390/electronics14010069.
  \url{https://www.mdpi.com/2079-9292/14/1/69}.
\newblock \textit{Evaluiert ML-Algorithmen unter verschiedenen Graden
  von Klassenungleichgewicht in IDS-Datensätzen. Beantwortet die Frage:
  Welche messbaren Konsequenzen hat Klassenungleichgewicht auf die
  Modellleistung, und warum ist Accuracy allein als Metrik ungenügend?
  Direkte Grundlage für Teilfrage 1 der eigenen Arbeit.}

\bibitem{ogunseyi2026xai_review}
T.~B. Ogunseyi, G.~Thiyagarajan, H.~He, V.~Bist, and Z.~Du.
\newblock Performance Analysis of Explainable Deep Learning-Based
  Intrusion Detection Systems for IoT Networks: A Systematic Review.
\newblock \emph{Sensors}, 26(2):363, 2026. doi:10.3390/s26020363.
  \url{https://www.mdpi.com/1424-8220/26/2/363}.
\newblock \textit{PRISMA-Review von 129 Studien (2018--2025). Belegt,
  dass SHAP in 98 Prozent der XAI-IDS-Studien eingesetzt wird ohne
  Genauigkeitseinbußen. Beantwortet die Frage: Ist die
  SHAP-Integration wissenschaftlich begründet und verbreitet in der
  IDS-Literatur?}

\bibitem{mohale2025xai}
V.~Z. Mohale and I.~C. Obagbuwa.
\newblock Evaluating Machine Learning-Based Intrusion Detection Systems
  with Explainable AI: Enhancing Transparency and Interpretability.
\newblock \emph{Frontiers in Computer Science}, 7:1520741, 2025.
  doi:10.3389/fcomp.2025.1520741.
  \url{https://www.frontiersin.org/journals/computer-science/articles/10.3389/fcomp.2025.1520741/full}.
\newblock \textit{Direkte methodische Vorlage für die SHAP-Integration.
  Zeigt konkret, wie globale und lokale SHAP-Analysen in einem
  IDS-Klassifikationsexperiment umgesetzt werden. Beantwortet die Frage:
  Wie sieht eine wissenschaftlich fundierte XAI-Analyse in einem
  IDS-Experiment aus?}

\bibitem{hermosilla2025xai_forensic}
P.~Hermosilla, S.~Berr{\'i}os, and H.~Allende-Cid.
\newblock Explainable AI for Forensic Analysis: A Comparative Study of
  SHAP and LIME in Intrusion Detection Models.
\newblock \emph{Applied Sciences}, 15(13):7329, 2025.
  doi:10.3390/app15137329.
  \url{https://www.mdpi.com/2076-3417/15/13/7329}.
\newblock \textit{Vergleicht SHAP und LIME auf XGBoost und TabNet im
  Kontext der Intrusion Detection. Führt konkrete Metriken zur Bewertung
  der Erklärungsqualität ein: Fidelität, Konsistenz und
  Jaccard-Ähnlichkeit. Beantwortet die Frage: Wie lässt sich die
  Qualität von SHAP-Erklärungen quantitativ und systematisch bewerten?
  Direkte Grundlage für die semantische Validierung in Teilfrage 3.}

\end{thebibliography}

\end{document}